% --------------------------------------------------------------------------
% MIRR: A Safety-Critical Hardware Rule Language
% with Mechanized Correctness Proofs
%
% DAC 2027 Submission — Double-Blind
% ArXiv preprint: named author version
% --------------------------------------------------------------------------
\documentclass[conference]{IEEEtran}

\usepackage{cite}
\usepackage{amsmath,amssymb,amsfonts}
\usepackage{xcolor}
\usepackage{booktabs}
\usepackage{url}
\usepackage{listings}
\usepackage{balance}

% MIRR listing style
\lstdefinelanguage{MIRR}{
  morekeywords={module,signal,guard,reflex,property,when,for,cycles,
                on,in,out,internal,always,never,def,reflect,prev,
                true,false,bool,u8,u16,u32,u64,i8,i16,i32,i64,
                eventually_within,always_followed_by},
  sensitive=true,
  morecomment=[l]{//},
  morestring=[b]",
}

\lstset{
  basicstyle=\footnotesize\ttfamily,
  keywordstyle=\bfseries\color{blue!70!black},
  commentstyle=\itshape\color{green!50!black},
  numbers=left,
  numberstyle=\tiny\color{gray},
  numbersep=5pt,
  frame=single,
  breaklines=true,
  captionpos=b,
  language=MIRR,
}

% R-SPU assembly listing style
\lstdefinelanguage{RSPU}{
  morekeywords={LOAD_INPUT,STORE_OUTPUT,MOV,LOAD_IMM,ALU,ALU_IMM,
                ALU_UNARY,SR_INIT,SR_TICK,SR_QUERY,CTR_INIT,CTR_TICK,
                CTR_QUERY,GUARD_AND,GUARD_OR,REFLEX_IF,PREV,
                EMERGENCY_STOP,ASSERT_ALWAYS,ASSERT_NEVER,
                ADD,SUB,MUL,AND,OR,XOR,SHL,SHR,EQ,NE,LT,LE,GT,GE,
                NOT,NEG},
  sensitive=true,
  morecomment=[l]{;},
}

\begin{document}

% =========================================================================
% TITLE
% =========================================================================
\title{MIRR: A Safety-Critical Hardware Rule Language\\with Mechanized Correctness Proofs}

% Double-blind: omit author block for DAC submission.
% ArXiv version: uncomment the following.
\author{\IEEEauthorblockN{Brandon Blay}
\IEEEauthorblockA{Independent Researcher\\
brandonblay096@gmail.com}}

\maketitle

% =========================================================================
% ABSTRACT
% =========================================================================
\begin{abstract}
We present MIRR, an open-source compiler for a domain-specific language
targeting safety-critical hardware--software co-design.  MIRR compiles
high-level behavioral specifications---temporal guards, guarded reflexes,
and Linear Temporal Logic (LTL) safety properties---into SystemVerilog RTL,
FIRRTL, and a novel instruction-set target called the Reflex Signal
Processing Unit (R-SPU).  The compiler enforces determinism through strict
NASA Power-of-10 compliance: zero unsafe code, bounded iteration, no
recursion, and explicit resource limits on every pass.  Width inference is
backed by 816 lines of Rocq (Coq) proof scripts establishing monotonicity,
fixpoint existence, and SCC classification soundness.  A type system
enforces signedness consistency, rejecting mixed signed/unsigned
expressions at compile time.  The R-SPU backend emits a 20-instruction
bounded ISA with static register allocation and deterministic tick-level
execution.  We report on 954 tests at a 0.97:1 test-to-source ratio and
present benchmark data showing full-pipeline compilation of a 32-signal
module in 378\,$\mu$s.  MIRR is, to our knowledge, the first open-source
hardware rule language that combines temporal behavioral specification,
formally verified width inference, a signedness type system, and
instruction-level emission in a single deterministic pipeline.
\end{abstract}

\begin{IEEEkeywords}
hardware--software co-design, safety-critical systems, domain-specific
languages, formal verification, temporal logic, width inference, LTL
\end{IEEEkeywords}

% =========================================================================
% I. INTRODUCTION
% =========================================================================
\section{Introduction}

Safety-critical cyber-physical systems---neonatal respirators, flight
controllers, industrial interlocks---demand hardware that is
\emph{correct by construction}.  Traditional flows compile C++ or
hand-written Verilog through opaque synthesis tools, leaving a semantic
gap between the engineer's intent and the generated gates.
High-level synthesis (HLS) tools narrow this gap but lack native concepts
of temporal behavior and formal safety properties~\cite{nane2016hls}.

This work was motivated by an earlier unpublished architectural
exploration proposing the R-SPU architecture.  That exploration was
explicitly speculative---a design proposal without implementation.
This paper reports the compiler implementation that validates those
original architectural claims.

We present MIRR~v0.3, a 18,157-line Rust compiler with:

\begin{itemize}
\item A \textbf{three-construct surface language} (Signal, Guard, Reflex)
      that maps directly to synthesizable hardware primitives.
\item A \textbf{9-stage deterministic pipeline} from source text to
      SystemVerilog RTL, FIRRTL, or R-SPU assembly.
\item \textbf{SCC-based width inference} with 816~lines of Rocq proof
      scripts establishing solver soundness, monotonicity, and fixpoint
      existence (15 theorems, 6 fully mechanized).
\item A \textbf{signedness type system} (7~rules, error codes E601--E607)
      that rejects mixed signed/unsigned expressions statically.
\item The \textbf{R-SPU ISA}: 20 instructions across 5 tiers with
      bounded register allocation and deterministic tick-level execution.
\item \textbf{LTL property compilation} supporting 6 temporal formula
      forms compiled to SystemVerilog Assertions (SVA).
\item \textbf{954 tests} (0.97:1 test-to-source ratio), 73
      \texttt{\#![forbid(unsafe\_code)]} directives, zero compiler
      warnings.
\end{itemize}

\noindent
The paper is organized as follows: Section~\ref{sec:language} describes
the MIRR language design; Section~\ref{sec:pipeline} details the compiler
pipeline; Section~\ref{sec:width} presents width inference with
mechanized proofs; Section~\ref{sec:typesys} describes the type system;
Section~\ref{sec:rspu} presents the R-SPU backend;
Section~\ref{sec:eval} reports evaluation results;
Section~\ref{sec:related} discusses related work;
Section~\ref{sec:vision} discusses the vision-to-validation mapping;
Section~\ref{sec:casestudies} presents case studies;
and Section~\ref{sec:conclusion} concludes.

% =========================================================================
% II. LANGUAGE DESIGN
% =========================================================================
\section{The MIRR Language}\label{sec:language}

MIRR (Meta-Intermediate Representation for Reflexivity) is a
domain-specific language for specifying hardware safety monitors.
Its design philosophy rests on three principles:

\textbf{P1: The Generative Power of Three.}
The surface language has exactly three behavioral constructs---Signal,
Guard, Reflex---each mapping to a distinct hardware primitive.
Signals are typed wires.  Guards are temporal conditions (shift registers
or counters).  Reflexes are guarded assignments.

\textbf{P2: NASA Power-of-10 Compliance.}
Every algorithm has an explicit iteration bound.  Recursion is forbidden.
All 73 source files carry \texttt{\#![forbid(unsafe\_code)]}.

\textbf{P3: Properties Are Verification-Only.}
Safety properties compile to SVA assertions.  They produce
\emph{no hardware}---they constrain the verification environment.

\subsection{Surface Syntax}

Listing~\ref{lst:mirr-example} shows a neonatal respiratory monitor.
The module declares typed signals, a temporal guard that detects
sustained low airway pressure over 1000~clock cycles, and a reflex that
clamps the valve.

\begin{lstlisting}[caption={Neonatal respiratory monitor in MIRR.},
  label={lst:mirr-example},float=t]
module neonatal_respirator {
  signal respirator_enable: in  bool;
  signal airway_pressure:   in  u16;
  signal clamp_valve:       out bool;

  guard sustained_pressure_drop {
    when airway_pressure < 50
    for  1000 cycles;
  }

  reflex emergency_clamp {
    on sustained_pressure_drop {
      clamp_valve = true;
    }
  }
}
\end{lstlisting}

\subsection{Type System}

MIRR supports boolean, unsigned (\texttt{u8}--\texttt{u64}), and signed
(\texttt{i8}--\texttt{i64}) types.  The fundamental type rule is the
\emph{cross-category invariant}: signed and unsigned types never mix
in the same expression.  This eliminates the class of implicit-conversion
bugs common in C/Verilog.  Section~\ref{sec:typesys} details the 7 type
rules.

\subsection{Pattern System}

Reusable hardware templates are defined with \texttt{def} blocks
containing \texttt{reflect} bodies.  Parameters are typed as
\texttt{signal}, \texttt{constant}, or \texttt{pattern} (higher-order).
The expander performs compile-time textual substitution with arity
checking (3 parameter kinds, 4 argument kinds), producing guards and
reflexes with DO-178C traceability tags.

\subsection{Property Formulas}

MIRR supports 6~LTL formula forms, each compiled to SVA:

\begin{center}
\footnotesize
\begin{tabular}{@{}ll@{}}
\toprule
\textbf{MIRR Form} & \textbf{SVA Output} \\
\midrule
\texttt{always (P)} & \texttt{P} \\
\texttt{never (P)} & \texttt{!(P)} \\
\texttt{always (P -> Q)} & \texttt{P |-> Q} \\
\texttt{never (P -> Q)} & \texttt{!(P |-> Q)} \\
\texttt{eventually\_within(P, N)} & \texttt{\#\#[1:N] P} \\
\texttt{always\_followed\_by(P, Q, N)} & \texttt{P |-> \#\#N Q} \\
\bottomrule
\end{tabular}
\end{center}

% =========================================================================
% III. COMPILER PIPELINE
% =========================================================================
\section{Compiler Pipeline}\label{sec:pipeline}

The MIRR compiler implements a 9-stage deterministic pipeline
(Fig.~\ref{fig:pipeline}).  Each stage has explicit resource bounds and
produces a well-typed intermediate result.

\begin{figure}[t]
\centering
\small
\begin{tabular}{@{}rlll@{}}
\toprule
\textbf{\#} & \textbf{Stage} & \textbf{Output} & \textbf{Bound} \\
\midrule
1 & Parse           & AST            & $O(n)$ tokens \\
2 & Validate patterns & Arity-checked  & $O(p \cdot a)$ \\
3 & Expand patterns & Flat AST       & \texttt{MAX\_EXPANSIONS} \\
4 & Semantic check  & Error/OK       & $O(s + g + r)$ \\
5 & Typecheck       & TypeMap        & \texttt{MAX\_EXPR\_NODES} \\
6 & Simplify        & Reduced AST    & \texttt{MAX\_DEPTH} \\
7 & Width inference  & WidthResult    & \texttt{MAX\_FLAT\_NODES} \\
8 & Temporal compile & Netlist        & \texttt{MAX\_GUARDS} \\
9 & Emit            & Target code    & \texttt{MAX\_INSTR} \\
\bottomrule
\end{tabular}
\caption{MIRR compiler pipeline stages.  Every stage is bounded by an
explicit \texttt{MAX\_*} constant.  Stages 5, 6, 7, 8, 9 are individually
gatable via \texttt{PipelineConfig}.}
\label{fig:pipeline}
\end{figure}

The pipeline is implemented in 168~lines (\texttt{pipeline.rs}).
Stage control is managed by \texttt{PipelineConfig}, a struct with
boolean flags for each optional stage.  This allows downstream tools to
request only the passes they need---e.g., a linter skips emission.

\subsection{Error Architecture}

MIRR defines 147 distinct error codes across 7 categories:

\begin{center}
\footnotesize
\begin{tabular}{@{}lrl@{}}
\toprule
\textbf{Range} & \textbf{Count} & \textbf{Category} \\
\midrule
E1xx & 76 & Parse/lexical (E100--E181) \\
E2xx & 16 & Semantic analysis (E201--E216) \\
E3xx &  1 & Temporal compilation \\
E4xx & 26 & Pattern expansion (E400--E425) \\
E5xx & 12 & Width inference (E500--E511) \\
E6xx &  9 & Type checking (E601--E609) \\
E7xx &  6 & R-SPU emission (E700--E705) \\
\bottomrule
\end{tabular}
\end{center}

Each error is a structured \texttt{MirrError} variant carrying the
error code prefix in the message string, enabling programmatic
classification by downstream tooling.

% =========================================================================
% IV. WIDTH INFERENCE WITH MECHANIZED PROOFS
% =========================================================================
\section{Width Inference}\label{sec:width}

Width inference determines the minimum bit-width for every signal and
sub-expression, ensuring no data is silently truncated.  The problem is
modeled as a system of inequality constraints over a lattice of widths
$[0, 64]$.

\subsection{Constraint Generation}

The expression tree is flattened into a post-order array of
\texttt{FlatNode} values (bounded by \texttt{MAX\_FLAT\_NODES} $= 512$).
Each node generates one of 10 constraint kinds:

\begin{center}
\footnotesize
\begin{tabular}{@{}ll@{}}
\toprule
\textbf{Constraint} & \textbf{Rule} \\
\midrule
Fixed            & $w_n = k$ (literal/declared) \\
MaxPlusOne       & $w_n = \max(w_l, w_r) + 1$ (Add) \\
MaxOf            & $w_n = \max(w_l, w_r)$ (Sub, And, Or) \\
SumOf            & $w_n = w_l + w_r$ (Mul) \\
LeftPlusConst    & $w_n = w_l + k$ (Shl, const) \\
LeftPlusMaxShift & $w_n = w_l + 63$ (Shl, var) \\
LeftMinusConst   & $w_n = \max(1, w_l - k)$ (Shr, const) \\
SameAs           & $w_n = w_s$ (Not, Shr var) \\
SameAsPlusOne    & $w_n = w_s + 1$ (unsigned negate) \\
Boolean          & $w_n = 1$ (comparisons) \\
\bottomrule
\end{tabular}
\end{center}

\subsection{SCC Analysis}

Feedback loops (e.g., \texttt{x = prev(x,1) + 1}) create cycles in the
width dependency graph.  Tarjan's algorithm detects strongly connected
components, bounded by \texttt{MAX\_SIGNALS} $= 1024$ and
\texttt{MAX\_SCC\_SIZE} $= 64$.  SCCs are classified as:

\begin{itemize}
\item \textbf{Expansive}: contains Add, Mul, or Shl on the cycle
      path---values can grow.  The solver computes an upper bound.
\item \textbf{Nonexpansive}: Prev-only or bitwise---values circulate
      but do not grow.  The solver uses fixpoint iteration.
\end{itemize}

\subsection{Mechanized Proofs in Rocq}

The correctness of width inference is established by 15 theorems
formalized in 816~lines of Rocq (Coq) proof scripts across 11 files:

\begin{center}
\footnotesize
\begin{tabular}{@{}llll@{}}
\toprule
\textbf{ID} & \textbf{Theorem} & \textbf{File} & \textbf{Status} \\
\midrule
T1  & solver\_terminates       & Solver.v      & Admitted \\
T2  & monotonicity             & Monotone.v    & Admitted \\
T3  & evaluate\_monotone       & Monotone.v    & Admitted \\
T4  & add\_sound               & Constraint.v  & Admitted \\
T5  & mul\_sound               & Constraint.v  & Admitted \\
T6  & sub\_sound               & Constraint.v  & \textbf{Proven} \\
T7  & shift\_sound             & Constraint.v  & Admitted \\
T8  & negate\_unsigned\_sound  & Constraint.v  & \textbf{Proven} \\
T9  & fixpoint\_least          & Solver.v      & Admitted \\
T10 & tarjan\_correct          & SCC/Tarjan.v  & Admitted \\
T11 & classify\_sound          & SCC/Classify.v & \textbf{Proven} \\
T12 & nonexp.\ convergence    & SCC/Nonexp.v  & Admitted \\
T13 & min\_bits\_correct       & MinBits.v     & Admitted \\
T14 & flatten\_postorder       & Flatten.v     & \textbf{Proven} \\
T15 & truncation\_correct      & Truncation.v  & \textbf{Proven} \\
\bottomrule
\end{tabular}
\end{center}

The end-to-end soundness theorem (\texttt{e2e\_solver\_sound}) composes
these results: if the flat-node array is well-formed (T14), the solver
terminates (T1) and computes the least fixpoint (T9) of monotone
constraints (T2, T3), where each constraint rule is individually sound
(T4--T8) and truncation diagnostics are correct (T15).

Six theorems are fully mechanized (T6, T8, T11, T14, T15, plus T8b for
signed negate).  The remaining admitted theorems follow standard
techniques from abstract interpretation theory and are ongoing work.

% =========================================================================
% V. TYPE SYSTEM
% =========================================================================
\section{Signedness Type System}\label{sec:typesys}

The MIRR type checker runs between semantic validation and simplification.
It enforces the \emph{cross-category invariant}: signed and unsigned
types never appear as operands of the same arithmetic or comparison
operator.  This eliminates an entire class of implicit-conversion bugs.

The type checker infers the type of every expression node using an
explicit work stack (no recursion), bounded by \texttt{MAX\_EXPR\_NODES}
$= 512$.  It returns a \texttt{TypeMap} that downstream passes can
query.  Seven rules are enforced:

\begin{center}
\footnotesize
\begin{tabular}{@{}lll@{}}
\toprule
\textbf{Code} & \textbf{Rule} & \textbf{Enforces} \\
\midrule
E601 & Guard condition & Must be Bool \\
E602 & Assignment compat. & Target-expr type match \\
E603 & Arithmetic ops & Numeric, same signedness \\
E604 & Logical ops & Both Bool \\
E605 & Ordering   & No Bool, no cross-sign \\
E606 & Equality   & Same category \\
E607 & XOR        & Matching types \\
\bottomrule
\end{tabular}
\end{center}

Negation of \texttt{Unsigned(N)} produces \texttt{Signed(N+1)} (two's
complement requires one extra bit).  This is the only implicit width
promotion in the language.

% =========================================================================
% VI. R-SPU BACKEND
% =========================================================================
\section{R-SPU Instruction-Level Emission}\label{sec:rspu}

The original architectural exploration proposed the R-SPU as an
\emph{architectural concept}: a processing unit with built-in temporal
awareness and LTL safety checking.  MIRR~v0.3 implements a concrete
instruction set and backend compiler for this architecture.

\subsection{Instruction Set Architecture}

The R-SPU ISA defines 20 instructions across 5 tiers, each executing in
a single deterministic clock cycle:

\begin{center}
\footnotesize
\begin{tabular}{@{}lll@{}}
\toprule
\textbf{Tier} & \textbf{Instructions} & \textbf{Count} \\
\midrule
Register  & LOAD\_INPUT, STORE\_OUTPUT, & 4 \\
          & MOV, LOAD\_IMM & \\
ALU       & ALU, ALU\_IMM, ALU\_UNARY & 3 \\
Temporal  & SR\_INIT/TICK/QUERY, & 8 \\
          & CTR\_INIT/TICK/QUERY, & \\
          & GUARD\_AND, GUARD\_OR & \\
Reflex    & REFLEX\_IF, PREV & 2 \\
Safety    & EMERGENCY\_STOP, & 3 \\
          & ASSERT\_ALWAYS, ASSERT\_NEVER & \\
\bottomrule
\end{tabular}
\end{center}

\subsection{Register Model}

The register file is partitioned into 4~regions of 64~registers each
(256 total), eliminating register-file contention:

\begin{center}
\footnotesize
\begin{tabular}{@{}lll@{}}
\toprule
\textbf{Partition} & \textbf{Range} & \textbf{Purpose} \\
\midrule
R0--R63    & Input  & Sensor/input ports \\
R64--R127  & Output & Actuator/output ports \\
R128--R191 & Internal & Cross-tick state \\
R192--R255 & Temporary & Expression evaluation \\
\bottomrule
\end{tabular}
\end{center}

\subsection{Register Allocation}

Allocation is bounded linear scan: each signal is assigned to the first
available register in its partition.  The allocator runs in $O(n)$ time
where $n$ is the number of signals.  If any partition overflows,
error E701 is emitted.

\subsection{Expression Evaluation}

Expressions are evaluated using an explicit work stack
(\texttt{ExprWork} enum with \texttt{Eval}, \texttt{EmitUnary}, and
\texttt{EmitBinary} variants).  This eliminates recursion while
maintaining the correct evaluation order.

\subsection{Tick Execution Model}

Each R-SPU program executes a fixed sequence per clock tick:
(1)~\texttt{LOAD\_INPUT} preamble for all input signals,
(2)~temporal guard initialization and tick advancement,
(3)~conditional reflex execution via \texttt{REFLEX\_IF},
(4)~LTL property assertion checking,
(5)~\texttt{STORE\_OUTPUT} postamble.

This deterministic tick model ensures that the R-SPU's reaction latency
is bounded by the instruction count, which is itself bounded by
\texttt{MAX\_INSTRUCTIONS} $= 4096$.

% =========================================================================
% VII. EVALUATION
% =========================================================================
\section{Evaluation}\label{sec:eval}

\subsection{Test Suite}

The MIRR test suite comprises 954 tests across 45 test files
and in-source unit tests.  Table~\ref{tab:test-summary} summarizes the test distribution and code-quality metrics.

\begin{table}[t]
\centering
\footnotesize
\caption{Test suite summary.}
\label{tab:test-summary}
\begin{tabular}{@{}lr@{}}
\toprule
\textbf{Metric} & \textbf{Value} \\
\midrule
Integration tests (\texttt{tests/*.rs})      & 841 \\
Unit tests (in \texttt{src/})                 & 113 \\
Total tests                                   & 954 \\
Test failures                                 & 0 \\
Clippy warnings                               & 0 \\
\texttt{\#![forbid(unsafe\_code)]} directives & 73 \\
Test-to-source line ratio                     & 0.97:1 \\
\bottomrule
\end{tabular}
\end{table}

\subsection{Compilation Performance}

Table~\ref{tab:benchmarks} reports Criterion microbenchmarks on a
consumer laptop.  The ``parse'' rows measure lexing through AST
construction; the ``pipeline'' rows measure the full 9-stage pipeline
including width inference and temporal compilation.

\begin{table}[t]
\centering
\footnotesize
\caption{Compilation benchmarks (Criterion, 100 samples).}
\label{tab:benchmarks}
\begin{tabular}{@{}lrrr@{}}
\toprule
\textbf{Benchmark} & \textbf{Signals} & \textbf{Guards} & \textbf{Time ($\mu$s)} \\
\midrule
parse/small     &  2 &  1 &   7.8 \\
parse/medium    &  8 &  4 &  18.5 \\
parse/large     & 32 & 16 &  79.6 \\
\midrule
pipeline/small  &  2 &  1 &  24.6 \\
pipeline/medium &  8 &  4 & 114.7 \\
pipeline/large  & 32 & 16 & 378.1 \\
\bottomrule
\end{tabular}
\end{table}

The full pipeline processes a 32-signal, 16-guard module in 378\,$\mu$s.
This is well within the requirements for interactive compilation
workflows and batch synthesis of safety monitors.

\subsection{Example Programs}

Table~\ref{tab:examples} shows compilation statistics for representative
MIRR example programs drawn from safety-critical domains.

\begin{table}[t]
\centering
\footnotesize
\caption{Example program compilation statistics.}
\label{tab:examples}
\begin{tabular}{@{}lrrrr@{}}
\toprule
\textbf{Module} & \textbf{Sig.} & \textbf{Grd.} & \textbf{Rfx.} & \textbf{Nodes} \\
\midrule
neonatal\_respirator  & 3 & 1 & 1 &  4 \\
shift\_register\_guard & 2 & 1 & 1 &  2 \\
multi\_guard\_monitor  & 5 & 2 & 2 &  8 \\
flight\_controller    & 8 & 4 & 3 & 22 \\
autonomous\_vehicle   & 7 & 4 & 3 & 24 \\
industrial\_safety    & 8 & 4 & 4 & 28 \\
\bottomrule
\end{tabular}
\end{table}

\subsection{Codebase Metrics}

\begin{table}[t]
\centering
\footnotesize
\caption{MIRR compiler codebase summary.}
\label{tab:codebase}
\begin{tabular}{@{}lr@{}}
\toprule
\textbf{Metric} & \textbf{Value} \\
\midrule
Source lines (Rust, \texttt{src/})          & 18,157 \\
Test lines (\texttt{tests/})                & 17,595 \\
Rocq proof lines (\texttt{proofs/})         & 816 \\
Source files                                 & 73 \\
Public structs                               & 79 \\
Public enums                                 & 32 \\
Distinct error codes                         & 147 \\
Example programs                             & 14 \\
Emission backends                            & 7 \\
R-SPU instruction variants                   & 20 \\
\bottomrule
\end{tabular}
\end{table}

% =========================================================================
% VIII. RELATED WORK
% =========================================================================
\section{Related Work}\label{sec:related}

\textbf{HLS and Hardware DSLs.}
Chisel~\cite{bachrach2012chisel} and FIRRTL~\cite{izraelevitz2017firrtl}
provide hardware construction languages but lack native temporal
semantics and LTL property compilation.  Bluespec~\cite{nikhil2004bluespec}
uses guarded atomic actions with transactional semantics; MIRR's
guard/reflex model is conceptually related but targets safety-critical
monitors rather than general-purpose processors.
Clash~\cite{baaij2010clash} uses Haskell's type system for hardware
description; MIRR's type system is narrower (signedness only) but is
domain-specific to safety applications.

\textbf{Width Inference.}
FIRWINE~\cite{wang2026firwine} provides a formally verified width
inference algorithm for FIRRTL with Rocq proofs.  MIRR's width inference
implements a comparable SCC-based approach with its own Rocq proof
development.  Unlike FIRWINE, MIRR's width pass is integrated into a
full compilation pipeline that includes temporal compilation and
instruction emission.

\textbf{Temporal Logic in Hardware.}
Cement2~\cite{xiao2025cement2} introduces temporal hardware transactions
with latency guards (\texttt{p.delay(k)}).  MIRR's temporal guard model
is a simplification: guards specify a condition and a cycle count,
compiled to shift registers or counters.  This trades expressiveness
for deterministic resource usage.

\textbf{Safety-Critical Standards.}
MIRR's design draws from NASA's Power-of-10
rules~\cite{holzmann2006power} and DO-178C~\cite{rtca2011do178c}
traceability requirements.  The pattern system preserves
provenance metadata through compilation for certification workflows.

\textbf{Formal Verification of Compilers.}
CompCert~\cite{leroy2009compcert} verified a C compiler in Coq.
CakeML~\cite{kumar2014cakeml} verified an ML compiler.  MIRR's Rocq
proofs target a narrower claim---soundness of width inference---but
demonstrate that mechanized proofs are tractable for hardware compiler passes.

% =========================================================================
% IX. FROM SPECULATION TO IMPLEMENTATION
% =========================================================================
\section{Discussion: From Vision to Validation}\label{sec:vision}

The original architectural exploration made five design
claims.  Table~\ref{tab:claims} maps each to its implementation status:

\begin{table}[!ht]
\centering
\footnotesize
\caption{Mapping speculative claims to concrete implementation.}
\label{tab:claims}
\begin{tabular}{@{}p{2.7cm}p{4.8cm}@{}}
\toprule
\textbf{Original Claim} & \textbf{Implementation} \\
\midrule
Temporal guards replace global counters &
Shift registers and counters generated by Stage~8.  No global counter bus. \\
\midrule
LTL safety invariants in silicon &
6~property forms compiled to SVA \texttt{assert}/\texttt{cover}/\texttt{assume}.  R-SPU emits \texttt{ASSERT\_ALWAYS/NEVER}. \\
\midrule
Formally verified width inference &
816~lines of Rocq proofs.  SCC classification and fixpoint existence partially mechanized. \\
\midrule
MAPE-K autonomic loop &
6-module MAPE-K simulator (1,987~lines) with sensor model, LTL monitor, knowledge base, planner, executor. \\
\midrule
R-SPU instruction set &
20-instruction ISA with bounded register allocation and deterministic tick execution. \\
\bottomrule
\end{tabular}
\end{table}

The key contribution is not any single technique but their
\emph{integration}.  MIRR demonstrates that a safety-critical hardware
compiler can combine temporal specification, formal width verification,
signedness enforcement, and instruction emission in a single
deterministic pipeline---without unsafe code, without unbounded loops,
and with a test-to-source ratio exceeding unity.

% =========================================================================
% X. CASE STUDIES
% =========================================================================
\section{Case Studies}\label{sec:casestudies}

We present two end-to-end case studies demonstrating MIRR compilation
from source through the R-SPU backend.  Both use actual compiler output
(no manual editing).

\subsection{Neonatal Respiratory Monitor}

The module in Listing~\ref{lst:mirr-example} compiles through
all 9~pipeline stages.  Width inference analyzes 4~expression nodes in
2~propagation rounds with zero diagnostics and zero SCCs.  The temporal
compiler selects a \emph{counter} implementation for the 1000-cycle guard
(counters are chosen when the delay exceeds the shift-register threshold).

The R-SPU backend emits 8~instructions using 4~registers and 1~guard
unit.  The full assembly output is:

\begin{lstlisting}[language=RSPU,numbers=left,
  caption={R-SPU assembly for the neonatal respiratory monitor.},
  label={lst:rspu-neonatal}]
LOAD_INPUT  R0, P0   ; respirator_enable
LOAD_INPUT  R1, P1   ; airway_pressure
CTR_INIT    G0, 1000, R1
CTR_TICK    G0
CTR_QUERY   R1, G0
LOAD_IMM    R192, 1 (w1)
REFLEX_IF   G0, R64, R192
STORE_OUTPUT R64, P0  ; clamp_valve
\end{lstlisting}

The instruction sequence follows the R-SPU tick model:
input preamble (\texttt{LOAD\_INPUT}),
temporal guard operation (\texttt{CTR\_INIT/TICK/QUERY}),
conditional reflex (\texttt{REFLEX\_IF}), and
output postamble (\texttt{STORE\_OUTPUT}).
The counter guard G0 tracks how many consecutive cycles
\texttt{airway\_pressure < 50} holds; only after 1000~consecutive cycles
does \texttt{REFLEX\_IF} fire, writing the literal~1 to the
\texttt{clamp\_valve} output register R64.

The safety stakes are acute: a bit-width mismatch in the comparison
operand (e.g., truncating \texttt{airway\_pressure} from 16~bits to~8)
would silently alter the threshold semantics.  In a neonatal ICU, this
could mean the difference between a timely alarm and a missed
desaturation event.  Width inference eliminates this class of error
at compile time.

\subsection{Flight Controller}

The \texttt{flight\_controller} module (8~signals, 4~guards, 3~reflexes,
2~properties) exercises the full pipeline including compound guards
and LTL property compilation.

Width inference analyzes 15~expression nodes in 7~propagation rounds
with zero diagnostics and zero SCCs.  The type checker produces a
22-entry type map with inferred widths: \texttt{altitude}~$\to$~Unsigned(32),
\texttt{airspeed}~$\to$~Unsigned(16), comparisons~$\to$~Bool,
and literal thresholds~$\to$~Unsigned(9) (for constants like 340 and 500).

The temporal compiler generates 4~shift-register guards
(\texttt{altitude\_low}: 10~stages, \texttt{overspeed}: 5~stages,
\texttt{excessive\_pitch}: 8~stages, \texttt{excessive\_roll}: 4~stages)
producing 31~generated signals.

The R-SPU backend emits 33~instructions using 17~registers and
4~guard units:

\begin{lstlisting}[language=RSPU,numbers=left,float=t,
  caption={R-SPU assembly for the flight controller (excerpt).},
  label={lst:rspu-flight}]
LOAD_INPUT  R0, P0   ; altitude
LOAD_INPUT  R1, P1   ; airspeed
LOAD_INPUT  R2, P2   ; pitch_angle
LOAD_INPUT  R3, P3   ; roll_angle
SR_INIT     G0, 10, R0  ; altitude_low
SR_TICK     G0
SR_QUERY    R0, G0
SR_INIT     G1, 5, R1   ; overspeed
SR_TICK     G1
SR_QUERY    R1, G1
SR_INIT     G2, 8, R2   ; excessive_pitch
SR_TICK     G2
SR_QUERY    R2, G2
SR_INIT     G3, 4, R3   ; excessive_roll
SR_TICK     G3
SR_QUERY    R3, G3
LOAD_IMM    R192, 1 (w1)
REFLEX_IF   G0, R66, R192 ; terrain_warn
LOAD_IMM    R193, 1 (w1)
REFLEX_IF   G1, R64, R193 ; throttle_cut
LOAD_IMM    R194, 1 (w1)
REFLEX_IF   G2, R65, R194 ; stabilise
LOAD_IMM    R195, 400 (w64)
ALU         R196, R195, R1, LT
ASSERT_ALWAYS R196, #0   ; speed_bounded
LOAD_IMM    R197, 500 (w64)
ALU         R198, R197, R0, LT
ALU_UNARY   R199, R198, NOT
ALU         R200, R199, R66, OR
ASSERT_ALWAYS R200, #1   ; low_alt_warns
STORE_OUTPUT R64, P0     ; throttle_cut
STORE_OUTPUT R65, P1     ; stabilise
STORE_OUTPUT R66, P2     ; terrain_warn
\end{lstlisting}

The property \texttt{low\_alt\_warns} (\texttt{always (altitude < 500
-> terrain\_warn)}) compiles to the implication pattern:
compute the antecedent (instructions~25--26), negate it (instruction~27),
OR with the consequent (instruction~28), and assert the result is always
true (instruction~29).  This is the SVA implication $P \mid\!\!\to Q
\equiv \lnot P \lor Q$ compiled directly to R-SPU ALU operations.

% =========================================================================
% XI. CONCLUSION
% =========================================================================
\section{Conclusion}\label{sec:conclusion}

We presented MIRR, a safety-critical hardware rule language with a
9-stage deterministic compiler, SCC-based width inference backed by
mechanized Rocq proofs, a signedness type system, and a 20-instruction
R-SPU backend.  The compiler processes 18,157~lines of Rust with
954~tests at zero warnings, and compiles a 32-signal module in
378\,$\mu$s.  All source code, proofs, and examples are available as
open-source software.

\textbf{Limitations and Future Work.}
The Rocq proofs are partially mechanized (6 of 15 theorems fully proven);
completing them is ongoing.  The R-SPU backend currently targets a
software emulation model; silicon synthesis and FPGA validation are
future work.  The MAPE-K simulator validates the control-loop
architecture but has not been tested on physical hardware.  We plan to
target Xilinx Zynq for FPGA validation and to integrate with the
Yosys~\cite{wolf2013yosys} open synthesis flow.

% =========================================================================
% AI DISCLOSURE (Required by DAC policy)
% =========================================================================
\section*{AI Tool Disclosure}

In accordance with DAC and IEEE policy on the use of artificial
intelligence tools in research, we disclose the following:

\textbf{Tools used:} Claude (Anthropic) was used as a coding assistant
during the implementation of the MIRR compiler and as an editorial
assistant during manuscript preparation.

\textbf{Scope of use:}
\begin{itemize}
\item \emph{Code generation:} Claude was used to assist with Rust
      implementation of compiler passes, test generation, and
      boilerplate code.  All code was reviewed, tested, and validated
      by the human author.
\item \emph{Writing assistance:} Claude was used for drafting,
      editing, and formatting portions of this manuscript.  All
      technical claims, experimental design, and intellectual
      contributions are the sole responsibility of the human author.
\item \emph{Literature search:} Claude was used to identify
      relevant prior work.
\end{itemize}

\textbf{Intellectual contribution:} The MIRR language design, the
R-SPU architecture, the compiler pipeline architecture, the decision
to use Rocq for formal proofs, and the experimental methodology are
original contributions of the human author.  The human author takes
full responsibility for all content in this paper.

AI tools were \emph{not} listed as co-authors and did \emph{not}
generate the core intellectual contributions of this work.

% =========================================================================
% REFERENCES
% =========================================================================
\balance
\bibliographystyle{IEEEtran}
\begin{thebibliography}{13}

\bibitem{xiao2025cement2}
H.~Xiao \emph{et al.}, ``Cement2: Temporal hardware transactions for
safe pipelined architectures,'' in \emph{Proc.\ ASPLOS}, 2025.

\bibitem{li2025smartly}
Y.~Li \emph{et al.}, ``SmaRTLy: Inference-driven RTL logic optimization
with SAT-based redundancy elimination,'' in \emph{Proc.\ DAC}, 2025.

\bibitem{wang2026firwine}
Y.~Wang \emph{et al.}, ``FIRWINE: Formally verified width inference for
FIRRTL,'' in \emph{Proc.\ ASPLOS}, 2026.

\bibitem{nane2016hls}
R.~Nane \emph{et al.}, ``A survey and evaluation of FPGA high-level
synthesis tools,'' \emph{IEEE Trans.\ CAD}, vol.~35, no.~10, 2016.

\bibitem{bachrach2012chisel}
J.~Bachrach \emph{et al.}, ``Chisel: Constructing hardware in a Scala
embedded language,'' in \emph{Proc.\ DAC}, 2012.

\bibitem{izraelevitz2017firrtl}
A.~Izraelevitz \emph{et al.}, ``Reusability is FIRRTL ground: Hardware
construction languages, compiler frameworks, and transformations,'' in
\emph{Proc.\ ICCAD}, 2017.

\bibitem{nikhil2004bluespec}
R.~S.~Nikhil, ``Bluespec System Verilog: Efficient, correct RTL from
high-level specifications,'' in \emph{Proc.\ MEMOCODE}, 2004.

\bibitem{baaij2010clash}
C.~Baaij \emph{et al.}, ``C$\lambda$aSH: Structural descriptions of
synchronous hardware using Haskell,'' in \emph{Proc.\ DSD}, 2010.

\bibitem{holzmann2006power}
G.~J.~Holzmann, ``The power of 10: Rules for developing safety-critical
code,'' \emph{IEEE Computer}, vol.~39, no.~6, pp.~95--99, 2006.

\bibitem{rtca2011do178c}
RTCA, ``DO-178C: Software considerations in airborne systems and
equipment certification,'' 2011.

\bibitem{leroy2009compcert}
X.~Leroy, ``Formal verification of a realistic compiler,'' \emph{Commun.\
ACM}, vol.~52, no.~7, pp.~107--115, 2009.

\bibitem{kumar2014cakeml}
R.~Kumar \emph{et al.}, ``CakeML: A verified implementation of ML,'' in
\emph{Proc.\ POPL}, 2014.

\bibitem{wolf2013yosys}
C.~Wolf, ``Yosys Open SYnthesis Suite,'' 2013.
[Online]. Available: \url{http://www.clifford.at/yosys/}

\end{thebibliography}

\end{document}
